A robótica de serviço tem experimentado um crescimento significativo, impulsionada pela demanda por automação em tarefas domésticas e cotidianas. No entanto, a massificação de robôs autônomos capazes de operar em ambientes residenciais enfrenta desafios técnicos e econômicos substanciais. O primeiro desafio reside na necessidade de navegação robusta em cenários dinâmicos — ambientes não estruturados onde a presença de pessoas, animais e a constante movimentação de mobiliário exigem algoritmos de localização resilientes \cite{abati2023slamvisualem}. O segundo desafio, e foco central deste trabalho, refere-se às restrições de \emph{hardware} impostas pela necessidade de baixo custo. Plataformas computacionais acessíveis e de baixo consumo energético frequentemente carecem dos recursos necessários para executar algoritmos de percepção de última geração em tempo real \cite{balado2025systematicliteraturereview}.

Neste contexto, abordagens de \emph{Visual SLAM} (vSLAM) surgem como uma alternativa promissora. O vSLAM consiste em um conjunto de algoritmos que capacita um dispositivo a construir o mapa de um ambiente desconhecido enquanto, simultaneamente, estima sua própria localização dentro dele, utilizando câmeras como fonte primária de percepção \cite{guapacha2017localizacaomapeamentosimultaneos}. Enquanto abordagens baseadas em LiDAR (\emph{Light Detection and Ranging}) com varredura de 360 graus são comuns na indústria pela sua precisão, esses sensores apresentam custo elevado e difícil aquisição no mercado nacional. Em contrapartida, o uso de câmeras RGB-D — capazes de capturar tanto a imagem colorida (RGB) quanto a informação de profundidade (Depth) — destaca-se por ser uma solução economicamente viável para mapeamento e localização em tempo real \cite{jin2019visualslamrgbd, lemos2019plataformaroboticamovel}.

Entretanto, a implementação de vSLAM em um computador de placa única, como o Raspberry Pi \cite{raspberrypi2019raspberrypi}, define o principal desafio desta pesquisa. O vSLAM é conhecido por exigir significativamente mais recursos de CPU do que o SLAM baseado em LiDAR, devido ao processamento intensivo de imagens. A literatura demonstra que algoritmos de ponta, como o ORB-SLAM3 \cite{campos2021orbslam3accurateopensource}, dificilmente atingem desempenho em tempo real em processadores limitados, levando a indústria a adotar plataformas mais onerosas com aceleração de GPU, como a família NVIDIA Jetson \cite{canovas2021onboarddynamicrgbd}. A escolha da plataforma Raspberry Pi para este trabalho justifica-se pela sua ampla adoção na comunidade acadêmica, baixo custo de aquisição, vasta documentação disponível e suporte nativo a sistemas operacionais baseados em Linux, essenciais para a execução do ROS. Esta escolha exige, portanto, uma análise cuidadosa do impacto do processamento de imagens e da gestão de memória.

Este trabalho parte da motivação de investigar e desenvolver uma arquitetura de navegação autônoma que priorize o baixo custo e a eficiência. A pesquisa adota a metodologia \emph{Design Science Research} (DSR) para conduzir a construção de um artefato tecnológico funcional \cite{hevner2004designscienceinformation, dresch2015distinctiveanalysiscase}. O artefato proposto consiste em um robô móvel construído sobre o chassi de um aspirador comercial reutilizado, visando a redução de custos mecânicos. A arquitetura integra um Raspberry Pi (como unidade de processamento de alto nível) a uma câmera RGB-D (como sensor de percepção), empregando um algoritmo de vSLAM para a navegação. Para o controle de baixo nível, utiliza-se um microcontrolador Arduino Nano \cite{arduino2024arduinonano} conectado a uma Ponte H para o acionamento dos motores. Este microcontrolador atua como uma ponte de comunicação entre o sistema operacional e os atuadores, além de permitir a integração e fusão com sensores adicionais para aprimorar a precisão da odometria e a robustez do mapeamento em situações onde a câmera RGB-D apresente limitações.

Para mitigar os riscos inerentes ao desenvolvimento de \emph{hardware} físico e permitir a validação da arquitetura em cenários diversos e controlados, este trabalho também adota uma estratégia de prototipagem simulada. Utilizando o simulador Gazebo \cite{koenig2004designuseparadigms}, será desenvolvido um modelo virtual do robô, permitindo testes de navegação e mapeamento em paralelo à montagem física. Esta abordagem dupla não apenas enriquece a análise de desempenho através da comparação entre o ambiente real e o simulado, mas também assegura a continuidade da pesquisa caso ocorram impedimentos críticos na implementação do \emph{hardware}, permitindo que a avaliação se concentre na validação algorítmica no ambiente virtual.

O sistema sob estudo deve ser orquestrado pelo \emph{Robot Operating System} (ROS), que provê a infraestrutura de comunicação necessária para integrar o fluxo de dados intenso do sensor visual com os comandos de atuação dos motores e a leitura dos demais elementos sensores \cite{quigley2009rosopensource}. O foco central do estudo é validar a viabilidade técnica desta arquitetura para a tarefa de navegação e mapeamento em ambientes domésticos desconhecidos, estabelecendo uma base robusta para trabalhos futuros focados em aplicações de serviço específicas, como a aspiração autônoma da residência.

A análise de desempenho é um componente crítico desta investigação. Busca-se caracterizar o comportamento do algoritmo vSLAM no \emph{hardware} limitado, monitorando métricas de consumo de CPU e uso de memória RAM. Com o foco nas restrições de processamento, especialmente no modo \emph{onboard}, o estudo avalia se a taxa de quadros por segundo (FPS, do inglês \textit{frames per second}) obtida no Raspberry Pi é compatível com uma navegação autônoma considerada segura. Neste trabalho, entende-se por \emph{navegação segura} o deslocamento sem colisões, mantendo distância mínima de obstáculos e respeitando limites de velocidade/rotação, além de uma taxa aceitável de falhas (por exemplo, perda de pose) durante a operação \cite{canovas2021onboarddynamicrgbd, nagarajan2024multiobjectiveoptimizationrtabmap, zheng2016multiframegraphmatching}.

\section{Objetivos}

Esta seção apresenta o objetivo geral e os objetivos específicos definidos para o desenvolvimento deste trabalho, detalhando as metas a serem alcançadas para a validação da proposta.


O objetivo geral deste trabalho é analisar a viabilidade técnica e o desempenho de uma arquitetura de navegação autônoma baseada em \emph{Visual SLAM} (vSLAM) aplicada a um robô de serviço com \emph{hardware} de recursos limitados. A proposta consiste em integrar um sensor RGB-D (Kinect) \cite{microsoft2012kinect_v1} a uma plataforma embarcada (Raspberry Pi) e validar a solução através de uma abordagem híbrida, que contempla tanto a simulação quanto a implementação física, focando na otimização do custo computacional e na robustez operacional em cenários residenciais.


Para alcançar o objetivo geral, foram definidos os seguintes objetivos específicos:
\begin{enumerate}
    \item Fundamentar teoricamente o funcionamento dos algoritmos e bibliotecas de \emph{Visual SLAM} (vSLAM) e definir o protocolo de avaliação experimental, levantando o estado da arte em navegação robótica de baixo custo;
    \item Modelar e simular o ambiente operacional e a cinemática do robô no \emph{software} Gazebo/ROS, permitindo a validação preliminar da integração e da lógica de navegação antes da implementação física;
    \item Integrar a arquitetura de \emph{hardware} do protótipo, incluindo sensores RGB-D (Kinect) e a camada de interface com hardware (Arduino Nano e atuadores);
    \item Integrar, configurar e testar a pilha de navegação baseada em vSLAM, empregando sensores RGB-D e, quando aplicável, fusão de sensores, de modo a mitigar erros de localização e aumentar a robustez do sistema;
    \item Avaliar o desempenho computacional da solução proposta, considerando métricas como consumo de CPU/RAM e taxa de quadros por segundo (FPS) durante a execução das tarefas de mapeamento e navegação.
\end{enumerate}


Considerando o exposto, a investigação será norteada pela seguinte questão de pesquisa:
\begin{center}
    \textit{É tecnicamente viável implementar um sistema de navegação autônoma robusto baseado em vSLAM com sensores de baixo custo (RGB-D) em uma plataforma embarcada limitada como o Raspberry Pi?}
\end{center}

\section{Organização do Trabalho}

Por fim, este documento está estruturado da seguinte forma: O Capítulo \ref{cap-fundamentacao-teorica} apresenta a fundamentação teórica sobre SLAM, vSLAM em \emph{hardware} de recursos limitados e as ferramentas de \emph{software} utilizadas. O Capítulo \ref{cap-metodologia} detalha a metodologia de pesquisa (DSR) e as etapas de construção do artefato. O Capítulo \ref{cap-trabalhos-relacionados} apresenta trabalhos correlatos que desenvolveram plataformas similares. O Capítulo \ref{cap-desenvolvimento} descreve o desenvolvimento do artefato nesta etapa, detalhando a arquitetura física do protótipo, a integração dos sensores e a validação preliminar da comunicação. Por fim, o Capítulo \ref{cap-conclusao} apresenta as considerações parciais, discutindo o cumprimento do cronograma, as limitações encontradas e o planejamento detalhado para a execução dos experimentos no Trabalho de Conclusão de Curso II.